% arara: xelatex
\documentclass[12pt]{article}

\usepackage{physics}
\usepackage{microtype}
\usepackage{array}
\usepackage{amsmath, amsfonts, amssymb}
\usepackage[top=2cm, left=1.6cm, right=1.6cm, bottom=2cm]{geometry}
\usepackage{booktabs}
\usepackage{enumitem}
\usepackage{fontspec}
\usepackage{polyglossia}
\usepackage{hyperref}

\setmainlanguage{russian}
\setotherlanguages{english}
\setmainfont{Arial}

\hypersetup{colorlinks=true,linkcolor=blue!60!black,urlcolor=blue!60!black}

\DeclareMathOperator{\Cov}{\mathbb{C}ov}
\DeclareMathOperator{\Corr}{\mathbb{C}orr}
\DeclareMathOperator{\Var}{\mathbb{V}ar}
\let\P\relax
\DeclareMathOperator{\P}{\mathbb{P}}
\DeclareMathOperator{\E}{\mathbb{E}}

\newcommand \e{\varepsilon}
\newcommand \cN{\mathcal{N}}
\newcommand{\cF}{\mathcal{F}}
\newcommand{\step}[1]{\medskip\noindent\textbf{#1}\par\nobreak\vspace{2pt}}

\setlength{\parindent}{0pt}
\setlength{\parskip}{6pt}

\title{Midterm 2026 --- решения}
\author{Временные ряды, ВШЭ, весна 2026}
\date{}

\begin{document}
\maketitle

\begin{center}
\textbf{Разбалловка.}\\[4pt]
\begin{tabular}{@{}l|ccc|c@{}}
\toprule
Задача & (a) & (b) & (c) & \textbf{Итого} \\
\midrule
1 --- ETS$(AAN)$                         & 5 & 5 & --- & 10 \\
2 --- GARCH$(1,1)$                       & 5 & 5 & 5   & 15 \\
3 --- ARIMA и прогноз                    & 5 & 5 & 5   & 15 \\
4 --- Лаговые полиномы и $MA(\infty)$    & 5 & 5 & 5   & 15 \\
5 --- Ковариационная функция             & 5 & 5 & --- & 10 \\
6 --- Нелинейный процесс                 & 5 & 5 & 5   & 15 \\
\midrule
\multicolumn{4}{r|}{\textbf{Всего за работу:}} & \(\mathbf{80}\) \\
\bottomrule
\end{tabular}
\end{center}

\section*{Задача 1. ETS$(AAN)$ \hfill [10 баллов]}

Рассматривается модель
\[
\begin{cases}
u_t \sim \cN(0;\sigma^2),\\
\ell_t = \ell_{t-1} + b_{t-1} + \alpha u_t,\\
b_t = b_{t-1} + \beta u_t,\\
y_t = \ell_{t-1} + b_{t-1} + u_t.
\end{cases}
\]

\step{(a) ARIMA-представление. \hfill [5 баллов]}
Начнём с уравнения наблюдения:
\[
y_t=\ell_{t-1}+b_{t-1}+u_t.
\]
Подставим в него уравнения состояния на один шаг назад:
\[
\ell_{t-1}=\ell_{t-2}+b_{t-2}+\alpha u_{t-1},
\qquad
b_{t-1}=b_{t-2}+\beta u_{t-1}.
\]
Тогда
\[
y_t=\ell_{t-2}+2b_{t-2}+(\alpha+\beta)u_{t-1}+u_t.
\]
При этом
\[
y_{t-1}=\ell_{t-2}+b_{t-2}+u_{t-1},
\]
поэтому
\[
y_t-y_{t-1}=b_{t-2}+(\alpha+\beta-1)u_{t-1}+u_t.
\]
Аналогично
\[
y_{t-1}-y_{t-2}=b_{t-3}+(\alpha+\beta-1)u_{t-2}+u_{t-1}.
\]
Из уравнения тренда
\[
b_{t-2}=b_{t-3}+\beta u_{t-2}.
\]
Вычитаем два соседних приращения:
\[
y_t-2y_{t-1}+y_{t-2}
=u_t+(\alpha+\beta-2)u_{t-1}+(1-\alpha)u_{t-2}.
\]
То есть
\[
(1-L)^2y_t=(1+(\alpha+\beta-2)L+(1-\alpha)L^2)u_t.
\]
Эквивалентно,
\[
\boxed{
y_t=2y_{t-1}-y_{t-2}+u_t+(\alpha+\beta-2)u_{t-1}+(1-\alpha)u_{t-2}.
}
\]
Значит, это модель $ARIMA(0,2,2)$ с параметрами
\[
\theta_1=\alpha+\beta-2,\qquad \theta_2=1-\alpha.
\]

\step{(b) Условные моменты. \hfill [5 баллов]}
По уравнениям состояния
\[
y_{t+1}=\ell_t+b_t+u_{t+1}.
\]
При условии на $\cF_t$ величины $\ell_t$ и $b_t$ известны, а $u_{t+1}$ независима от прошлого и имеет дисперсию $\sigma^2$. Поэтому
\[
\boxed{\E(y_{t+1}\mid\cF_t)=\ell_t+b_t,\qquad
\Var(y_{t+1}\mid\cF_t)=\sigma^2.}
\]

\section*{Задача 2. GARCH$(1,1)$ \hfill [15 баллов]}

Рассматривается модель
\[
r_t=\sigma_tz_t,\qquad z_t\sim\cN(0,1),
\]
\[
\sigma_t^2=1+0.25r_{t-1}^2+0.6\sigma_{t-1}^2.
\]

\step{(a) Безусловная дисперсия. \hfill [5 баллов]}
Так как
\[
0.25+0.6=0.85<1,
\]
существует стационарное решение GARCH-процесса с конечной второй моментной характеристикой. Далее используем это стационарное решение: именно поэтому математические ожидания в рекурсивном уравнении не зависят от $t$.

Обозначим
\[
v=\Var(r_t)=\E(r_t^2)=\E(\sigma_t^2).
\]
Тогда
\[
v=1+0.25v+0.6v,
\]
откуда
\[
0.15v=1,\qquad v=\frac{20}{3}.
\]
Следовательно,
\[
\boxed{\Var(r_t)=\frac{20}{3}.}
\]

\step{(b) Двухшаговый прогноз волатильности. \hfill [5 баллов]}
При $r_T=2$ и $\sigma_T^2=5$:
\[
\E(\sigma_{T+1}^2\mid\cF_T)
=1+0.25\cdot 2^2+0.6\cdot 5=5.
\]
Далее
\[
\E(r_{T+1}^2\mid\cF_T)=\E(\sigma_{T+1}^2\mid\cF_T)=5,
\]
поэтому
\[
\E(\sigma_{T+2}^2\mid\cF_T)
=1+0.25\cdot 5+0.6\cdot 5=5.25.
\]
Итак,
\[
\boxed{\E(\sigma_{T+1}^2\mid\cF_T)=5,\qquad
\E(\sigma_{T+2}^2\mid\cF_T)=5.25.}
\]

\step{(c) Долгосрочный предел. \hfill [5 баллов]}
На больших горизонтах прогноз условной дисперсии стационарного GARCH-процесса возвращается к безусловной дисперсии. В пункте (a) мы уже нашли
\[
v=\Var(r_t)=\E(r_t^2)=\E(\sigma_t^2)=\frac{20}{3}.
\]
Кроме того,
\[
r_{T+h}^2=\sigma_{T+h}^2z_{T+h}^2,\qquad \E(z_{T+h}^2)=1,
\]
поэтому прогноз $\E(r_{T+h}^2\mid\cF_T)$ имеет тот же долгосрочный предел, что и прогноз $\E(\sigma_{T+h}^2\mid\cF_T)$. Следовательно,
\[
\boxed{\lim_{h\to\infty}\E(r_{T+h}^2\mid\cF_T)=\frac{20}{3}.}
\]

\section*{Задача 3. ARIMA и прогноз \hfill [15 баллов]}

Рассматривается уравнение
\[
(1-0.7L)(1-L)y_t=3+(1+0.4L)\e_t,\qquad \e_t\sim WN(0,1).
\]

\step{(a) Классификация. \hfill [5 баллов]}
После одного взятия разности $w_t=\Delta y_t=(1-L)y_t$ получаем
\[
(1-0.7L)w_t=3+(1+0.4L)\e_t.
\]
Значит, у исходного процесса один единичный корень, AR-часть порядка 1 и MA-часть порядка 1:
\[
\boxed{y_t\sim ARIMA(1,1,1).}
\]

\step{(b) Первые значения ACF и PACF после взятия разности. \hfill [5 баллов]}
Для центрированного процесса $w_t-\E w_t$ имеем ARMA$(1,1)$:
\[
(1-0.7L)(w_t-\E w_t)=(1+0.4L)\e_t.
\]
Обозначим
\[
x_t=w_t-\E w_t.
\]
Тогда
\[
x_t=0.7x_{t-1}+\e_t+0.4\e_{t-1}.
\]
Найдём автоковариации напрямую. Для $k\ge 2$ умножаем уравнение на $x_{t-k}$ и берём ожидание:
\[
\gamma_k=0.7\gamma_{k-1}
+\Cov(\e_t,x_{t-k})+0.4\Cov(\e_{t-1},x_{t-k}).
\]
При $k\ge 2$ оба ковариационных слагаемых равны нулю, потому что $x_{t-k}$ зависит только от более старых шоков. Значит,
\[
\gamma_k=0.7\gamma_{k-1},\qquad k\ge 2.
\]
В частности,
\[
\rho_2=0.7\rho_1.
\]

Теперь найдём $\gamma_1$ и $\gamma_0$ из системы уравнений. Умножим
\[
x_t=0.7x_{t-1}+\e_t+0.4\e_{t-1}
\]
на $x_{t-1}$ и возьмём ожидание:
\[
\gamma_1=0.7\gamma_0+\Cov(\e_t,x_{t-1})+0.4\Cov(\e_{t-1},x_{t-1}).
\]
Здесь
\[
\Cov(\e_t,x_{t-1})=0,
\qquad
\Cov(\e_{t-1},x_{t-1})=\Cov(\e_{t-1},\e_{t-1})=1.
\]
Значит,
\[
\gamma_1=0.7\gamma_0+0.4.
\]

Теперь умножим исходное уравнение на $x_t$:
\[
\gamma_0=0.7\gamma_1+\Cov(\e_t,x_t)+0.4\Cov(\e_{t-1},x_t).
\]
Имеем
\[
\Cov(\e_t,x_t)=1,
\]
а также
\[
\Cov(\e_{t-1},x_t)
=0.7\Cov(\e_{t-1},x_{t-1})+\Cov(\e_{t-1},\e_t)
+0.4\Cov(\e_{t-1},\e_{t-1})
=0.7+0.4=1.1.
\]
Поэтому второе уравнение:
\[
\gamma_0=0.7\gamma_1+1+0.4\cdot 1.1
=0.7\gamma_1+1.44.
\]
Получили систему
\[
\begin{cases}
\gamma_1=0.7\gamma_0+0.4,\\
\gamma_0=0.7\gamma_1+1.44.
\end{cases}
\]
Решая её, получаем
\[
\gamma_0=\frac{172}{51},\qquad \gamma_1=\frac{704}{255}.
\]
Следовательно,
\[
\rho_1=\frac{\gamma_1}{\gamma_0}
=\frac{704/255}{172/51}
=\frac{176}{215}.
\]
И тогда
\[
\rho_2=0.7\rho_1=\frac{616}{1075}.
\]

В общем виде для ARMA$(1,1)$ с параметрами $\phi$ и $\theta$ получилась формула
\[
\rho_1=\frac{(\phi+\theta)(1+\phi\theta)}{1+2\phi\theta+\theta^2},
\qquad
\rho_k=\phi^{k-1}\rho_1,\quad k\ge 2.
\]

Теперь найдём первые два значения PACF через систему Юла--Волкера.

Для лага 1 система состоит из одного уравнения:
\[
\rho_1=\alpha_1.
\]
Значит,
\[
\alpha_1=\frac{176}{215}.
\]

Для лага 2 коэффициенты линейной проекции $x_t$ на $x_{t-1},x_{t-2}$ обозначим $\phi_{21}$ и $\phi_{22}$. Система Юла--Волкера:
\[
\begin{pmatrix}
1 & \rho_1\\
\rho_1 & 1
\end{pmatrix}
\begin{pmatrix}
\phi_{21}\\
\phi_{22}
\end{pmatrix}
=
\begin{pmatrix}
\rho_1\\
\rho_2
\end{pmatrix}.
\]
Второй PACF равен последнему коэффициенту этой системы:
\[
\alpha_2=\phi_{22}.
\]
Из системы получаем
\[
\phi_{22}
=\frac{\rho_2-\rho_1^2}{1-\rho_1^2}
=\frac{\frac{616}{1075}-\left(\frac{176}{215}\right)^2}
{1-\left(\frac{176}{215}\right)^2}
=-\frac{88}{299}\approx -0.294.
\]
Итак,
\[
\boxed{\rho_1=\frac{176}{215},\quad \rho_2=\frac{616}{1075},\quad
\alpha_1=\frac{176}{215},\quad \alpha_2=-\frac{88}{299}.}
\]

\step{(c) Долгосрочный прогноз разности. \hfill [5 баллов]}
Снова обозначим
\[
w_t=\Delta y_t=y_t-y_{t-1}.
\]
После взятия разности исходное уравнение становится
\[
(1-0.7L)w_t=3+(1+0.4L)\e_t,
\]
то есть
\[
w_t=3+0.7w_{t-1}+\e_t+0.4\e_{t-1}.
\]
У стационарного решения среднее постоянно: $\E(w_t)=\E(w_{t-1})=\mu_w$. Поскольку $\E(\e_t)=0$ и $\E(\e_{t-1})=0$, берём математическое ожидание:
\[
\mu_w=3+0.7\mu_w.
\]
Отсюда
\[
\mu_w=\frac{3}{1-0.7}=10.
\]

Для долгосрочного прогноза центрируем процесс:
\[
x_t=w_t-\mu_w.
\]
Тогда
\[
x_t=0.7x_{t-1}+\e_t+0.4\e_{t-1}.
\]
Так как $|0.7|<1$, влияние текущей информации на прогноз $x_{T+h}$ убывает как степени $0.7$ и стремится к нулю. Поэтому условный прогноз $w_{T+h}$ стремится к безусловному среднему:
\[
\E(w_{T+h}\mid\cF_T)\to \mu_w=10.
\]
Так как $w_{T+h}=y_{T+h}-y_{T+h-1}$, получаем
\[
\boxed{
\lim_{h\to\infty}\E(y_{T+h}-y_{T+h-1}\mid\cF_T)=10.
}
\]

\section*{Задача 4. Лаговые полиномы и $MA(\infty)$ \hfill [15 баллов]}

Рассматривается уравнение
\[
y_t-0.9y_{t-1}+0.2y_{t-2}=5+u_t-0.5u_{t-1},
\qquad u_t\sim WN(0,1).
\]

\step{(a) Разложение на множители. \hfill [5 баллов]}
Левая часть:
\[
1-0.9L+0.2L^2=(1-0.5L)(1-0.4L).
\]
Правая часть:
\[
1-0.5L.
\]
Итак,
\[
\boxed{(1-0.9L+0.2L^2)=(1-0.5L)(1-0.4L),\qquad
1-0.5L \text{ уже разложен.}}
\]

\step{(b) Более простое уравнение. \hfill [5 баллов]}
Сначала найдём среднее стационарного решения:
\[
(1-0.9+0.2)\mu=5,
\]
то есть
\[
0.3\mu=5,\qquad \mu=\frac{50}{3}.
\]
Тогда
\[
(1-0.5L)(1-0.4L)(y_t-\mu)=(1-0.5L)u_t.
\]
Для стационарных решений общий множитель $(1-0.5L)$ можно сократить:
\[
(1-0.4L)(y_t-\mu)=u_t.
\]
Эквивалентно,
\[
y_t=10+0.4y_{t-1}+u_t.
\]
Следовательно,
\[
\boxed{(1-0.4L)\left(y_t-\frac{50}{3}\right)=u_t}
\quad\text{или}\quad
\boxed{y_t=10+0.4y_{t-1}+u_t.}
\]

\step{(c) Явный вид $MA(\infty)$. \hfill [5 баллов]}
Так как $|0.4|<1$,
\[
y_t-\frac{50}{3}=(1-0.4L)^{-1}u_t
=\sum_{j=0}^{\infty}0.4^j u_{t-j}.
\]
Отсюда
\[
\boxed{
y_t=\frac{50}{3}+u_t+0.4u_{t-1}+0.4^2u_{t-2}+\cdots.
}
\]

\section*{Задача 5. Ковариационная функция \hfill [10 баллов]}

Для $j=1,2$ задано $\gamma^{(j)}_{-k}=\gamma^{(j)}_k$. Кроме того, при $k\ge 0$:
\[
\gamma^{(1)}_k=
\begin{cases}
1, & k=0,\\
0.8, & k=1,\\
0.2, & k=2,\\
0, & k\ge 3,
\end{cases}
\qquad
\gamma^{(2)}_k=
\begin{cases}
1, & k=0,\\
0, & k=1,\\
\frac12, & k=2,\\
0, & k\ge 3.
\end{cases}
\]

\step{(a) Проверка существования. \hfill [5 баллов]}
Для $\gamma^{(1)}$ рассмотрим ковариационную матрицу в моменты $t,t+1,t+2$:
\[
\Gamma_3=
\begin{pmatrix}
1 & 0.8 & 0.2\\
0.8 & 1 & 0.8\\
0.2 & 0.8 & 1
\end{pmatrix}.
\]
Её определитель равен
\[
\det\Gamma_3=-0.064<0.
\]
Ковариационная матрица обязана быть неотрицательно определённой, поэтому $\gamma^{(1)}$ не может быть ковариационной функцией стационарного процесса.

Для $\gamma^{(2)}$ приведём явный пример в пункте (b), значит эта функция является допустимой.
Итак,
\[
\boxed{\gamma^{(1)}\text{ не может быть ковариационной функцией,}\qquad
\gamma^{(2)}\text{ может.}}
\]

\step{(b) Явный пример. \hfill [5 баллов]}
Пусть
\[
x_t=\eta_t+\eta_{t-2},\qquad \eta_t\sim WN\left(0,\frac12\right).
\]
Тогда
\[
\Var(x_t)=\frac12+\frac12=1,
\]
\[
\Cov(x_t,x_{t-1})=0,
\]
а
\[
\Cov(x_t,x_{t-2})=\Var(\eta_{t-2})=\frac12.
\]
При $|k|\ge 3$ общих шумов нет, поэтому ковариации равны нулю. Следовательно,
\[
\boxed{x_t=\eta_t+\eta_{t-2},\quad \eta_t\sim WN\left(0,\frac12\right)}
\]
имеет ковариационную функцию $\gamma^{(2)}$.

\section*{Задача 6. Нелинейный процесс \hfill [15 баллов]}

Пусть $(\e_t)$ --- независимые случайные величины, принимающие значения $-1$ и $1$ с вероятностями $1/2$ и $1/2$. Рассмотрим процесс
\[
x_t=\e_t\e_{t-1}.
\]

\step{(a) Слабая стационарность. \hfill [5 баллов]}
Проверим только слабую стационарность. Математическое ожидание не зависит от $t$:
\[
\E(x_t)=\E(\e_t)\E(\e_{t-1})=0,
\]
поскольку $\E(\e_t)=0$. Дисперсия также постоянна:
\[
\Var(x_t)=\E(\e_t^2\e_{t-1}^2)=1.
\]
Для лага 1:
\[
\Cov(x_t,x_{t-1})
=\E(\e_t\e_{t-1}^2\e_{t-2})=0.
\]
Для $|k|\ge 2$ произведение $x_tx_{t-k}$ содержит независимые множители с нулевыми средними, поэтому
\[
\Cov(x_t,x_{t-k})=0.
\]
Ковариационная функция зависит только от лага:
\[
\gamma_0=1,\qquad \gamma_k=0,\quad k\ne 0.
\]
Следовательно,
\[
\boxed{(x_t)\text{ является слабо стационарным процессом.}}
\]

\step{(b) Белый шум. \hfill [5 баллов]}
Для лага 1:
\[
\Cov(x_t,x_{t-1})
=\E(\e_t\e_{t-1}^2\e_{t-2})=0.
\]
Для $|k|\ge 2$ произведение $x_tx_{t-k}$ тоже содержит независимые множители с нулевыми средними, поэтому
\[
\Cov(x_t,x_{t-k})=0,\qquad k\ne 0.
\]
Итак,
\[
\boxed{(x_t)\text{ является белым шумом.}}
\]

\step{(c) Независимость. \hfill [5 баллов]}
Проверим независимость последовательно. Сначала заметим, что каждое $x_t$ принимает значения $-1$ и $1$ с вероятностями $1/2$ и $1/2$.

Возьмём два соседних значения, например $x_t$ и $x_{t+1}$. Для любых $a,b\in\{-1,1\}$ условия
\[
x_t=a,\qquad x_{t+1}=b
\]
означают
\[
\e_t\e_{t-1}=a,\qquad \e_{t+1}\e_t=b.
\]
Можно свободно выбрать $\e_t$: два варианта. После этого $\e_{t-1}$ и $\e_{t+1}$ определяются однозначно. Значит, из $2^3$ равновероятных наборов $(\e_{t-1},\e_t,\e_{t+1})$ подходят два:
\[
\P(x_t=a,x_{t+1}=b)=\frac{2}{8}=\frac14
=\P(x_t=a)\P(x_{t+1}=b).
\]
То есть соседние $x_t$ и $x_{t+1}$ независимы.

Для нескольких подряд идущих величин рассуждение такое же. Пусть заданы
\[
a_1,\ldots,a_n\in\{-1,1\}.
\]
Тогда система
\[
x_1=a_1,\quad \ldots,\quad x_n=a_n
\]
эквивалентна
\[
\e_1\e_0=a_1,\quad \ldots,\quad \e_n\e_{n-1}=a_n.
\]
Свободно выбираем $\e_0$: два варианта. После этого все $\e_1,\ldots,\e_n$ определяются последовательно и однозначно. Поэтому
\[
\P(x_1=a_1,\ldots,x_n=a_n)=\frac{2}{2^{n+1}}=\frac{1}{2^n}.
\]
Но каждое $x_i$ равно $a_i$ с вероятностью $1/2$, так что
\[
\P(x_1=a_1,\ldots,x_n=a_n)
=\prod_{i=1}^n \P(x_i=a_i).
\]
Значит, величины $x_t$ независимы:
\[
\boxed{(x_t)\text{ состоит из независимых случайных величин.}}
\]

\end{document}
